\subsection{Semisimple Modules}

DEFINITION 1. --- A module is called semisimple if it is the direct sum of a family of simple modules. \(^{(1)}\) .

A multimodule is called semisimple if it is the direct sum of a family of simple multimodules (cf.~I, §4, No.~4, p.~37, Definition 7).

An A-module M is semisimple if and only if it is semisimple when viewed as a module over its ring of homotheties \(A_{M}\) .

Examples. --- 1) A module reduced to 0 and a simple module are semisimple modules.

\begin{enumerate}
\def\labelenumi{\arabic{enumi})}
\setcounter{enumi}{1}
\item
  If A is a field, then every A-module is semisimple by Theorem 1 of II, §7, No.~1, p.~292. This shows that, in general, a semisimple module decomposes in several ways into a direct sum of simple submodules (see, however, Corollary 2 of VIII, p.~68).
\item
  Let A be a principal ideal domain, and let P be a system of representatives consisting of irreducible elements of A (VII, §1, No.~3, p.~3). Let M be an A-module, and, for every \(\pi \in \mathrm{P}\) , let \(\mathrm{M}(\pi)\) be the set of \(x \in \mathrm{M}\) such that \(\pi x = 0\) . By VII, §2, No.~2, p.~9, the A-module M is semisimple if and only if it is the sum of the submodules \(\mathrm{M}(\pi)\) ; it is then the direct sum of these submodules. This example will be generalized further on (VIII, p.~65).
\end{enumerate}

Let \(A_{1}\) and \(A_{2}\) be algebras over a commutative ring K. In III, §4, No.~3, p.~466, we introduced the notion of left bimodule over the algebras \(A_{1}\) and \(A_{2}\) and showed that this notion is equivalent to that of a left module over the ring \(A_{1} \otimes_{K} A_{2}\) . We say that M is a simple (resp. semisimple, finitely generated) bimodule if it is a simple (resp. semisimple, finitely generated) module over the ring \(A_{1} \otimes_{K} A_{2}\) .

THEOREM 1. --- Let M be a module that is the (not necessarily direct) sum of a family \((\mathrm{S}_{i})_{i\in\mathrm{I}}\) of simple submodules, and let N be a submodule of M. There exists a subset J of I such that M is the direct sum of the family consisting of N and the modules \(S_{j}\) for j running through J.

Let \(\mathcal{S}\) be the set of subsets \(I'\) of \(I\) such that the sum of the family consisting of the modules \(N\) and \(S_i\) for \(i\) in \(I'\) is direct. The set \(\mathcal{S}\) is of finite character: a subset \(J\) of \(I\) belongs to \(\mathcal{S}\) if and only if the same holds for every finite subset of \(J\) . Hence, the set \(\mathcal{S}\) has a maximal element \(J\) (Set Theory, III, §4, No.~5, p.~171). Set \(N' = N + \sum_{j \in J} S_j\) . Let \(i\) be in \(I - J\) . Since \(J\) is maximal in \(\mathcal{S}\) , the set \(J \cup \{i\}\) does not belong to \(\mathcal{S}\) , so that \(S_i \cap N' \neq 0\) . Since \(S_i\) is a simple module, we have \(S_i \cap N' = S_i\) . We therefore have \(S_i \subset N'\) for every \(i \in I\) , whence \(N' = M\) . This completes the proof.

COROLLARY 1. --- Every module that is the sum of a family of simple modules is semisimple.

It suffices to apply Theorem 1 to the case \(\mathbf{N} = 0\) .

COROLLARY 2. --- A module M is semisimple if and only if every submodule of M is a direct factor.

The condition is necessary by Theorem 1.

Conversely, suppose that every submodule of M admits a supplement. Let \(M'\) be the sum of the simple submodules of M, and let \(M''\) be a supplement of \(M'\) in M. Suppose that we have \(M' \neq M\) and therefore \(M'' \neq 0\) . Let N be a nonzero monogenous submodule of \(M''\) . By Proposition 3 of VIII, p.~49, there exists a maximal submodule P of N. Let Q be a submodule supplementary to P in M. Then \(N \cap Q\) is a submodule of N, supplementary to P in N, hence isomorphic to N/P (II, §1, No.~9, p.~210, Proposition 13). Consequently, \(N \cap Q\) is a simple submodule of \(M''\) , contrary to the definition of \(M'\) .

We therefore have \(\mathrm{M}' = \mathrm{M}\) , and the module \(\mathrm{M}\) is semisimple by Corollary 1.

COROLLARY 3. --- Let M be a semisimple module and N a submodule of M. The modules N and M/N are semisimple. More precisely, if M is the direct sum of a family \((\mathrm{S}_i)_{i\in \mathrm{I}}\) of simple modules, then there exists a subset J of I such that M/N is isomorphic to \(\bigoplus_{j\in \mathrm{J}}\mathrm{S}_j\) and N to \(\bigoplus_{i\in \mathrm{I - J}}\mathrm{S}_i\) .

Choose J as in Theorem 1. The module \(N' = \bigoplus_{j \in J} S_j\) is supplementary to N in M; it is therefore isomorphic to M/N. Moreover, the submodules N and \(\bigoplus_{i \in I-J} S_i\) of M are both supplementary to \(N'\) and therefore isomorphic to \(M/N'\) .

COROLLARY 4. --- Let M be a semisimple module. Then M is simple if and only if the endomorphism ring E of M is a field.

If M is simple, then E is a field by the corollary of Proposition 2 of VIII, p.~47.

If E is a field, then the module M is indecomposable (VIII, p.~31, Proposition 4, a)). Since it is moreover semisimple, it is simple.

Remark. --- Let K be an algebraically closed commutative field and A a K-algebra. Let M be a semisimple A-module that is a finite-dimensional vector space over the field K. Then M is simple if and only if every endomorphism of the A-module M is of the form \(x \mapsto \alpha x\) with \(\alpha\) in K: this is necessary by Theorem 1 of VIII, p.~47, and sufficient by Corollary 4 above.

